La Terre (rayon de \qty{6.371e6}{m}) reçoit du Soleil une intensité de \qty{1361}{W/m^2} et son albédo est de \num{\v{al}}. Sa surface, de température moyenne \qty{\v{Ts}}{K}, rayonne pratiquement comme un corps noir dans l'infrarouge.

\begin{enumerate}[label=\alph*)]
    \item Déterminez la puissance solaire totale absorbée par la Terre.
    \item Déterminez la puissance totale rayonnée par la surface terrestre.
    \item Comparez les deux résultats. D'où vient la puissance supplémentaire que reçoit la surface pour se maintenir à \qty{\v{Ts}}{K} ?
\end{enumerate}
